\documentclass[a4paper,12pt]{article}

\input{../../preambolo.tex}

\usepackage[utf8]{inputenc}
\usepackage[T1]{fontenc}
\usepackage{listings}
\usepackage{xcolor}
\usepackage{geometry}
\geometry{margin=2.5cm}


\lstset{
	basicstyle=\ttfamily\small,
	breaklines=true,
	backgroundcolor=\color{gray!5},
	frame=single,
	keywordstyle=\color{blue},
	commentstyle=\color{teal},
	numberstyle=\tiny,
	numbers=left,
	xleftmargin=1em,
	showstringspaces=false
}

\title{Spiegazione del programma in Assembly x86}
\author{Liam Ferretti}
\date{}

\begin{document}
	
	\maketitle
	
	\section*{Struttura logica del programma}
	
	Il programma realizza la somma di due numeri naturali inseriti dall’utente.  
	L’operazione di somma viene scomposta in una serie di \textbf{operazioni elementari}, ognuna delle quali corrisponde a un blocco logico ben definito del programma.  
	Ogni blocco è identificato da una o più \textit{etichette} (\textit{label}) che fungono da punti di riferimento per i salti e i cicli di controllo.  
	
	L’organizzazione complessiva può essere descritta come segue:
	
	\begin{enumerate}
		\item \textbf{Sezione \texttt{.data}}:  
		Contiene tutti i dati costanti usati dal programma, come le stringhe di testo che vengono mostrate a schermo (ad esempio i messaggi che chiedono all’utente di inserire un numero).  
		Include anche le costanti che rappresentano la lunghezza di ciascun messaggio e il carattere di nuova linea.  
		
		\item \textbf{Sezione \texttt{.bss}}:  
		Serve per riservare spazio in memoria a variabili che verranno inizializzate in un secondo momento.  
		Qui troviamo le celle di memoria in cui verranno salvati i numeri inseriti dall’utente (\texttt{n1} e \texttt{n2}) e il \texttt{buffer} utilizzato per la conversione del risultato finale in ASCII, in modo da poterlo stampare.  
		
		\item \textbf{Label \texttt{\_start}}:  
		È il punto di ingresso del programma, ossia l’istruzione da cui il sistema operativo Linux inizia l’esecuzione. Da qui parte la sequenza di operazioni successive.  
		
		\item \textbf{Stampa e lettura del primo numero}:  
		Il programma mostra a schermo un messaggio con la system call \texttt{write} (numero 4) e poi attende l’input dell’utente usando la system call \texttt{read} (numero 3).  
		Il valore letto è una stringa di caratteri ASCII, cioè le cifre del numero digitate dall’utente, seguite da un carattere di newline.  
		
		\item \textbf{\texttt{.loop1} / \texttt{.done1}}:  
		Questa coppia di etichette identifica un ciclo che scorre la stringa contenente il primo numero e lo converte da caratteri ASCII a valore intero.  
		Ogni cifra viene isolata, trasformata nel corrispondente valore numerico e “accumulata” in \texttt{eax}, che diventa così il contenitore del numero convertito.  
		
		\item \textbf{Stampa e lettura del secondo numero}:  
		Il programma ripete la stessa sequenza di istruzioni per il secondo numero: prima stampa il messaggio \texttt{msg2}, poi legge l’input e lo memorizza in \texttt{n2}.  
		
		\item \textbf{\texttt{.loop2} / \texttt{.done2}}:  
		È la seconda routine di conversione ASCII → intero, del tutto analoga alla prima.  
		Alla fine di questo blocco, \texttt{eax} contiene il valore del secondo numero appena inserito.  
		
		\item \textbf{Somma}:  
		Una volta terminata la conversione, basta una sola istruzione \texttt{add eax, [n1]} per sommare il contenuto della cella \texttt{n1} (il primo numero) al valore già presente in \texttt{eax} (il secondo numero).  
		Il risultato della somma rimane memorizzato in \texttt{eax}.  
		
		\item \textbf{\texttt{.conv\_loop}}:  
		Questo ciclo svolge l’operazione inversa della conversione iniziale: trasforma il risultato numerico contenuto in \texttt{eax} in una sequenza di caratteri ASCII, memorizzandola nel \texttt{buffer}.  
		La divisione per 10 permette di estrarre le cifre una alla volta a partire da quella meno significativa, che vengono poi convertite in caratteri e salvate in memoria in ordine inverso.  
		
		\item \textbf{Stampa del risultato e terminazione}:  
		Infine, il programma mostra il messaggio “La somma è:” seguito dal numero convertito.  
		Dopo la stampa, viene eseguita una system call \texttt{exit} (numero 1) per terminare correttamente il programma e restituire il controllo al sistema operativo.
	\end{enumerate}
	
	\bigskip
	
	\noindent
	In sintesi, l’intero programma si può vedere come una catena di trasformazioni:  
	\[
	\text{Input ASCII} \;\rightarrow\; \text{Conversione in interi} \;\rightarrow\; \text{Somma} \;\rightarrow\; \text{Conversione in ASCII} \;\rightarrow\; \text{Output}
	\]
	Ogni etichetta rappresenta un tassello di questa catena, in cui le operazioni complesse sono costruite combinando passaggi elementari di lettura, aritmetica e scrittura gestiti tramite \texttt{syscall} e cicli controllati manualmente.
	
	
	\newpage
	
	\section*{Sezione \texttt{.data}}
	\begin{lstlisting}[language={[x86masm]Assembler}]
section .data
	msg1    db "Inserisci il primo numero (max 10 cifre): ", 0
	len1    equ $ - msg1
	msg2    db "Inserisci il secondo numero (max 10 cifre): ", 0
	len2    equ $ - msg2
	msgRes  db "La somma e': ", 0
	lenRes  equ $ - msgRes
	newline db 0Ah 
	\end{lstlisting}
	
	Questa sezione contiene i dati costanti del programma:
	\begin{itemize}
		\item Le stringhe da stampare sullo schermo.
		\item Le variabili \texttt{len1}, \texttt{len2}, \texttt{lenRes} che contengono la lunghezza delle rispettive stringhe, calcolata come \texttt{\$ - msgX}.
		\item \texttt{newline} è il carattere di nuova linea (codice ASCII 0Ah).
		\item lo zero dopo i db è il null terminator e serve a indicare la fine della stringa.
	
	\end{itemize}
	\section*{Sezione \texttt{.bss}}
	\begin{lstlisting}[language={[x86masm]Assembler}]
section .bss 
	buffer resd 10 
	n1 resd 1
	n2 resd 1
	\end{lstlisting}
	
	La sezione \texttt{.bss} serve per riservare spazio in memoria per dati variabili:
	\begin{itemize}
		\item \texttt{buffer} sarà usato per la conversione del risultato in ASCII.
		\item \texttt{n1} e \texttt{n2} conterranno i numeri interi letti da input.  
	\end{itemize}
	Il direttivo \texttt{resd} riserva "doublewords" (4 byte ciascuna), il \texttt{resw} riserva "words" (2 byte ciascuna), il \texttt{resb} riserva "byte" (1 byte ciascuna), il \texttt{resq} riserva "quadwords" ovvero 8 byte.
	
	\section*{Sezione \texttt{.text} e inizio del programma}
	\begin{lstlisting}[language={[x86masm]Assembler}]
section .text 
global _start

_start:
	\end{lstlisting}
	
	Qui inizia il codice eseguibile.  
	L’etichetta \texttt{\_start} indica il punto d’ingresso del programma (usato da Linux per sapere da dove iniziare l’esecuzione).
	
	\section*{Stampa del primo messaggio e lettura del primo numero}
	\begin{lstlisting}[language={[x86masm]Assembler}]
mov eax, 4
mov ebx, 1
mov ecx, msg1
mov edx, len1
int 0x80
	\end{lstlisting}
	
	Queste istruzioni invocano la \textbf{system call write} (numero 4) per stampare \texttt{msg1} sullo schermo.  
	I registri significano:
	\begin{itemize}
		\item \texttt{eax = 4} → numero della syscall \texttt{write}.
		\item \texttt{ebx = 1} → file descriptor dello standard output.
		\item \texttt{ecx = msg1} → indirizzo del messaggio.
		\item \texttt{edx = len1} → lunghezza della stringa.
	\end{itemize}
	
	\begin{lstlisting}[language={[x86masm]Assembler}]
mov eax, 3
mov ebx, 0
mov ecx, n1
mov edx, 10
int 0x80
	\end{lstlisting}
	
	Questa volta la syscall è \texttt{read}:
	\begin{itemize}
		\item \texttt{eax = 3} → numero della syscall \texttt{read}.
		\item \texttt{ebx = 0} → file descriptor dello standard input.
		\item \texttt{ecx = n1} → indirizzo su cui verranno salvati i dati.
		\item \texttt{edx = 10} → lunghezza dell'input.
	\end{itemize} 
	
	\newpage
	
	\section*{Conversione del primo numero da ASCII a intero}
	\begin{lstlisting}[language={[x86masm]Assembler}]
	xor eax, eax
	xor ebx, ebx
	mov esi, n1
.loop1:
	mov bl, [esi]
	cmp bl, 0x0A
	je .done1
	sub bl, '0'
	imul eax, eax, 10
	add eax, ebx
	inc esi
	jmp .loop1
	.done1:
	mov [n1], eax
	\end{lstlisting}
	
	Questo ciclo:
\begin{enumerate}
	\item Le prime due righe servono a resettare i valori in eax e ebx, mettendo in entrambi 0.
	\item muoviamo in esi l'indirizzo in memoria di n1, rendonolo quindi un puntatore che scorre i byte dell'input.
	\item .loop1, segno l'inizio del ciclo.
	\item \texttt{mov bl, [esi]}, carica in bl (gli ultimi 8 bit di ebx) il byte puntato da esi, infatti con [  ] si rappresenta il valore dell'indice in memoria, e contiene un carattere ASCII.
	\item \texttt{cmp bl, 10} confronta il valore di bl con 10, ovvero 0x0A in hex, che corrisponde al carattere newline ($\backslash$n).
	\item Ora \texttt{je .done1} nel caso in cui \texttt{cmp bl, 10} abbia restituito che bl è uguale a 10, allora salta a \texttt{.done1}
	\item Sottraggo da bl il carattere 0, \texttt{sub bl, '0'} dato che 0 in hex è 0x30, e i numeri da 0 a 9 sono rappresentati da 0x30 a 0x39, si ottiene quindi il numero intero.
	\item \texttt{imul eax, eax, 10} moltiplica eax per 10, e salva il risultato in eax, per spostare in avanti le cifre, per poter sommare la cifra delle unità
	\item \texttt{add eax, ebx} somma la cifra corrente, in bl, su eax, ma essendo una operazione a 32 bit, bisogna usare ebx.
	\item \texttt{inc esi} incremento il puntatore al byte successivo.
	\item \texttt{jmp .loop1} torno al punto 3.
	\item \texttt{.done1} etichetta di uscita dal ciclo 
	\item \texttt{mov [n1], eax}, scrive il valore intero del risultato in memoria all'indirizzo n1
\end{enumerate}
	
	\section*{Lettura e conversione del secondo numero}
	\begin{lstlisting}[language={[x86masm]Assembler}]
	mov eax, 4
	mov ebx, 1
	mov ecx, msg2
	mov edx, len2
	int 0x80
	
	mov eax, 3
	mov ebx, 0
	mov ecx, n2
	mov edx, 10
	int 0x80
	
	xor eax, eax
	xor ebx, ebx
	mov esi, n2
.loop2:
	mov bl, [esi]
	cmp bl, 0x0A
	je .done2
	sub bl, '0'
	imul eax, eax, 10
	add eax, ebx
	inc esi
	jmp .loop2
.done2:
	\end{lstlisting}
	
	È lo stesso procedimento del primo numero:  
	si legge la stringa da tastiera, poi si converte in valore numerico.
	
	\section*{Somma dei due numeri}
	\begin{lstlisting}[language={[x86masm]Assembler}]
add eax, [n1]
	\end{lstlisting}
	
	Dopo la conversione del secondo numero, \texttt{eax} lo contiene.  
	Qui si aggiunge \texttt{n1} (che contiene il primo numero).  
	Il risultato rimane in \texttt{eax}.
	
	\newpage
	
	\section*{Conversione del risultato in ASCII}
	\begin{lstlisting}[language={[x86masm]Assembler}]
	mov edi, buffer + 10
	mov ecx, 10
.conv_loop:
	dec edi
	xor edx, edx
	div ecx
	add dl, '0'
	mov [edi], dl
	test eax, eax
	jnz .conv_loop
	\end{lstlisting}
	
	Questo ciclo esegue la conversione da numero binario ad ASCII:
	\begin{enumerate}
		\item \texttt{mov edi,  buffer + 10}, metti nel registro edi, l'indirizzo alla fine del buffer, ovvero 10 byte.
		\item \texttt{mov ecx, 10}, carico in ecx il valore 10, per eseguire la divisione.
		\item \texttt{.conv\_loop:} etichetta per l'inizio del ciclo.
		\item \texttt{dec edi} sposto il puntatore di edi indietro di un byte verso sinistra
		\item azzero edx con xor
		\item \texttt{div ecx} divide il valore a 64 bit in edx:eax per 10 (ecx), e il quoziente finisce in eax, il resto in edx
		\item \texttt{add dl, '0'} converte la cifra in ASCII
		\item \texttt{mov [edi], dl} scrive la cifra ASCII in memoria puntata da edi
		\item \texttt{test eax, eax} verifica che il contenuto di eax sia 0 quindi abbiamo finito di dividere il numero.
		\item \texttt{jnz .conv\_loop} se eax è diverso da 0, quindi quando non abbiamo ancora diviso tutte le cifre.
	\end{enumerate}
	
	\newpage
	
	\section*{Stampa del risultato e terminazione}
	\begin{lstlisting}[language={[x86masm]Assembler}]
mov eax, 4
mov ebx, 1
mov ecx, msgRes
mov edx, lenRes
int 0x80

mov eax, 4
mov ebx, 1
mov ecx, edi
mov edx, buffer + 10
sub edx, edi
int 0x80

mov eax, 4
mov ebx, 1
mov ecx, newline
mov edx, 1
int 0x80

mov eax, 1
xor ebx, ebx
int 0x80
	\end{lstlisting}
	
	Infine:
	\begin{itemize}
		\item Stampa la stringa “La somma è:”.
		\item Stampa il numero convertito.
		\item Va a capo.
		\item Termina il programma con la syscall \texttt{exit (1)}. 
	\end{itemize}
	
\end{document}
